\begin{abstract}
We prove the ordinary two-point Chowla conjecture. For every fixed pair
of nonproportional affine forms, the Liouville correlation has a
power-of-logarithm saving at every cutoff, with an absolute exponent.
We also prove the binary corrected Elliott conjecture for ordinary
averages of complex multiplicative functions of modulus at most one,
under uniform nonpretentiousness of at least one original factor.
This qualitative conclusion holds in fixed residue classes and for
fixed nonproportional affine forms.
\end{abstract}

\section{Introduction}\label{sec:introduction}

For \(n\ge1\), let \(\Omega(n)\) be the number of prime factors of \(n\),
counted with multiplicity, and let \(\lambda(n)=(-1)^{\Omega(n)}\).
We study ordinary two-point averages: each integer up to the given
cutoff has weight one. Thus the correlation measures whether the two prime-factor counts
have the same parity more often than different parities on the full
initial interval.
The distinction from a logarithmic average is the distinction between
\[
 \frac1X\sum_{n\le X}a(n)
 \qquad\text{and}\qquad
 \frac1{\log X}\sum_{n\le X}\frac{a(n)}n.
\]
Logarithmic weighting combines information from many multiplicative
scales; an ordinary estimate must control the terminal scale itself.
Our quantitative conclusion concerns every fixed pair of
nonproportional affine forms.

\begin{theorem}[A logarithmic saving for affine Liouville correlations]
\label{thm:q-affine}
There is an absolute constant \(c>0\) such that, for every fixed choice
of integers \(a_1,a_2\ge1\) and \(b_1,b_2\ge0\) with
\(a_1b_2-a_2b_1\ne0\), there is a constant
\(C_{a_1,a_2,b_1,b_2}\) for which
\[
 \left|\sum_{1\le n\le X}
       \lambda(a_1n+b_1)\lambda(a_2n+b_2)\right|
 \le C_{a_1,a_2,b_1,b_2}\frac{X}{(\log X)^c}
 \qquad(X\ge3).
\]
The cutoff \(X\) may be any real number. No coprimality condition is
imposed on the coefficients.
\end{theorem}

For $a_1=a_2=1$ and $b_1=0$, $b_2=h>0$, this resolves the
ordinary two-point Chowla conjecture, with a logarithmic saving.
The exponent is independent of the forms. The implied constants need
not be effective, and no uniformity for coefficients growing with \(X\)
is asserted. The proof first obtains the same saving for
\(\lambda(n)\lambda(n+h)\) in each fixed residue class modulo each fixed
positive integer, including nonunit and zero classes. Complete
multiplicativity then gives the affine conclusion.

We also prove a qualitative statement for general multiplicative
functions. Let \(\mathbb D=\{z\in\C:|z|\le1\}\). A function
\(f:\N\to\mathbb D\) is multiplicative when
\(f(mn)=f(m)f(n)\) for coprime positive integers \(m,n\).
For a Dirichlet character \(\chi\), a real number \(t\), and \(X\ge2\),
define the nonnegative distance \(D\) by
\begin{equation}\label{eq:distance}
 D(f,\chi n^{it};X)^2
 =\sum_{p\le X}\frac{1-\operatorname{Re}
       (f(p)\overline{\chi(p)p^{it}})}p.
\end{equation}
All sums indexed by \(p\) are over primes. We call \(f\)
\emph{uniformly nonpretentious} if, for every fixed Dirichlet character,
\begin{equation}\label{eq:unp}
 \inf_{|t|\le N}D(f,\chi n^{it};N)\longrightarrow\infty
 \qquad(N\longrightarrow\infty).
\end{equation}
The prime cutoff and the height bound are the same \(N\).

\begin{theorem}[Binary corrected Elliott]\label{thm:elliott}
Let \(f_1,f_2:\N\to\mathbb D\) be multiplicative, and suppose that
at least one satisfies \eqref{eq:unp}. For every fixed pair of distinct
nonnegative integers \(h_1,h_2\),
\[
 \frac1N\sum_{n=1}^N f_1(n+h_1)f_2(n+h_2)\longrightarrow0
 \qquad(N\longrightarrow\infty,\ N\in\N).
\]
\end{theorem}

This proves the binary corrected Elliott conjecture for ordinary
averages. The functions may be complex, need not have modulus one,
and need not be completely multiplicative. There is no conjugation
in the displayed product; the conjugated version follows because
\eqref{eq:unp} is preserved by conjugation. We make no quantitative
rate claim for this class of functions.

\begin{corollary}[Affine correlations]\label{cor:affine}
Let \(f_1,f_2:\N\to\mathbb D\) be multiplicative, with at least one
satisfying \eqref{eq:unp}. For all fixed integers \(a_1,a_2\ge1\) and
\(b_1,b_2\ge0\) with \(a_1b_2-a_2b_1\ne0\),
\[
 \frac1N\sum_{n=1}^N f_1(a_1n+b_1)f_2(a_2n+b_2)
 \longrightarrow0.
\]
\end{corollary}

Every fixed progression restriction is also permitted in
\Cref{thm:elliott}; \Cref{prop:affine-progression} gives the precise
statement through all real cutoffs and for every residue class.
These conclusions use the hypothesis on an original factor \(f_j\).
The finite local changes needed for dilations and residue restrictions
preserve that hypothesis. In particular, Liouville satisfies the exact
condition \eqref{eq:unp}; a quantitative distance bound is proved in
\Cref{lem:affine-liouville-unp}.

